% ============================================================
%  Generic Essay Template
%  Compile: pdflatex essay.tex  (run twice)
% ============================================================
\documentclass[11pt,a4paper]{article}

% ---------- Packages ----------
\usepackage[T1]{fontenc}
\usepackage[utf8]{inputenc}
\usepackage{mathpazo}                 % Palatino font (remove for default)
\usepackage[margin=1in]{geometry}     % page margins
\usepackage{microtype}                % nicer spacing
\usepackage{setspace}                 % line spacing
\usepackage{enumitem}                 % list control
\usepackage{graphicx}                 % images
\usepackage{booktabs}                 % clean tables
\usepackage{csquotes}                 % \enquote{...}
\usepackage[hidelinks]{hyperref}      % clickable links

\onehalfspacing                       % or \singlespacing / \doublespacing
\setlength{\parskip}{0.5em}
\setlength{\parindent}{0pt}

% ---------- Placeholders (edit these) ----------
\newcommand{\EssayTitle}{Where Is the World Heading?}
\newcommand{\EssaySubtitle}{2026}
\newcommand{\EssayAuthor}{Ravikant Pothukuchi}
\newcommand{\EssayDate}{\today}

\title{\EssayTitle\\[0.3em]\large\itshape\EssaySubtitle}
\author{\EssayAuthor}
\date{\EssayDate}

% ============================================================
\begin{document}

\maketitle

NVIDIA is a prime example of a behemoth that can rise and fall. Its secrets
are out in the open now. That still leaves the real question open: where is
the world heading?

What I learnt from NVIDIA's story is that its bet was simple. GPUs could take
on processing that CPUs could not, and NVIDIA made that power programmable.
Now it is reaping the fruits. Its strength is seeing the future---next year,
the year after, and the year after that---and if it can bring everyone on
board, it wins.

\section*{The cosmic clock}

So where is the world heading? 2026 falls in the
second half of Brahma's life, in his 51st year. We are in the Sveta Varaha
Kalpa, the first day of that year. Within this day we are in the 7th of 14
Manvantaras, and in the 28th of its 71 Mahayugas. Within the Kali Yuga, this
is year 5128. Counted from the start of this kalpa, the year is
1,97,29,49,128---one billion, nine hundred seventy-two million, nine hundred
forty-nine thousand, one hundred twenty-eight years.

Where is the world heading? To year 1,97,29,49,129. And ultimately, towards
cosmic dissolution. At the end of the kalpa the physical, material world
dissolves. The jivas fall dormant, merging into Vishnu until the next cycle.
For the length of one kalpa the universe remains dissolved; then the dormant
karmas drive creation once again in a new kalpa.

Within this kalpa, manvantara, yuga and samvatsara, though, the world is
headed in one direction: towards the top of the food chain.

\section*{Everything is capital}

Movies, cars, houses, chips, AI, mail, computers, software, electronics,
satellites, aeroplanes, jets, missiles, drones, insurance products, banks,
plugins, browsers, droughts, food, animals, oil, countries, people, slaves,
weapons, guns, coins, gold, silver, concepts, ideas, yoga, fruits, prompts,
graphics, imagination, paintings, music, compact discs, Blu-rays, players,
vessels, stoves, air conditioners, water purifiers, plastics, chemicals,
medicines, survival kits, tents, tools, lamps, toilets, tubs, shower heads,
screws, nuts, glasses, LEDs, LCDs, rockets, telescopes, business logic, data,
privacy, cameras, printers, paper, news, articles, lands, islands, planets,
moons, suns, switches, adapters, phone cases, bags, guitars, helmets,
cushions, phones, mics, sofas, beds, travel plans, entertainment displays,
drivers, salaries, wires, copper, metals, books, memory.

That is the last twenty years in a nutshell, and we are heading towards more
and more of it. Everyone is a capitalist, capitalising on everything. From the
universe down to the microscopic, we are all searching for the next thing to
capitalise on. We do not ask whether it is good or bad, useful or destructive,
clean or polluting, honest or corrupt, beneficial or harmful---to children,
animals, plants, planets, moons, or ourselves. The question is far behind us,
the answer is far ahead, and we are stuck in between.

Every mind grabs on to anything it can. People are not just people; people
are their minds. And those minds are polluted, stimulated beyond any way back.

So where are we headed? Towards bigger, better, more efficient, more lustrous,
shinier, faster, stronger, more powerful, more dominant, ever-growing, sharper,
more rigid, lighter, easier, more stable, more capitalisable. Anything that
fits these adjectives has a future. We want to fly faster, but safer. We want
to strike with weapons without being struck ourselves. We want shields without
the burden of carrying them. We want infinite food without worrying about
starving the rest of the world. We want intelligence without worrying about
who dies for lack of water while it cools our data centres.

\section*{Better every year}

2026, 2027, 2028, 2029, 2030, 2031, 2032, 2033, 2034, 2035: each year has to
be better than the last. Anything that fits that equation survives the year.
The rest can rise from its ashes.

My identity as an engineer alone would limit what I see. So I identify first
as God, then as a son, a student, a husband, a father, an Indian, an engineer,
an IT engineer. Engineering comes near
the end. God, son and student come much earlier. That is who I am: a student
of this life and this universe, with the universe as my guru. From there, I
can see past the filter of bigger, better, faster, more efficient, more
powerful, more dominant, shinier, stronger, lighter, easier.

2027 and 2028 demand exactly that of me---of me as a human---and that is
what will carry this human forward:

\begin{itemize}[leftmargin=1.5em, itemsep=0.1em]
  \item the most efficient software;
  \item the most powerful software;
  \item the largest codebase, encompassing all of human knowledge;
  \item the largest memory for holding data;
  \item the smallest chip holding the most data;
  \item the fastest engine for computing;
  \item the easiest robot to configure;
  \item the most encompassing general-purpose platform.
\end{itemize}

\end{document}